\par\noindent\textbf{(1)} 取球心实验室坐标x=X+s cosα、y=y0−s sinα，故平动动能T\_cm=m[(Ẋ+ṡcosα)²+ṡ²sin²α]/2，楔块为M Ẋ²/2。纯滚动相对楔面令球自旋角速率ω=ṡ/R，旋转能κm ṡ²/2。总L=(M+m)Ẋ²/2+mẊṡcosα+m(1+κ)ṡ²/2+mg s sinα。
\par\noindent\textbf{(2)} X坐标无显含，给(M+m)Ẍ+mcosα s̈=0；s方程 mcosα Ẍ+m(1+κ)s̈=mg sinα。消元得 s̈=g sinα/[1+κ−mcos²α/(M+m)]。分母由(M+m)(1+κ)−mcos²α>0保证动力学正定，故球沿斜面向下加速。
\par\noindent\textbf{(3)} 楔块加速度 Ẍ=−[mcosα/(M+m)]s̈，方向与球沿斜面水平分量相反。水平总动量p\_X=(M+m)Ẋ+m ṡcosα守恒，这是楔和球内部接触力相互抵消的结果；不可假设球在实验室水平速度为零。
\par\noindent\textbf{(4)} M→∞时 s̈→g sinα/(1+κ)，回到固定斜面滚动加速度；κ→0给不旋转的点质量在移动斜面上的滑动极限。若无耗散，总能E=(M+m)Ẋ²/2+mẊṡcosα+m(1+κ)ṡ²/2−mg s sinα，沿方程求导为零。
